% !TEX root = main.tex

\section{Progettazione del dataset}
\label{sec:dataset}

\subsection{Fonte e universo}
I dati di mercato sono stati scaricati da Yahoo Finance tramite la libreria \texttt{yfinance} con l'opzione \texttt{auto\_adjust=True}. 
Questa opzione fornisce i dati già aggiustati per split e dividendi su base giornaliera in modo da evitare che il modello percepisca delle crisi quando in realtà è solo l'effetto causato da questi ultimi.
Nel dataset sono presenti titoli large-cap del mercato americano, cioè titoli che rappresentano un mercato molto liquido ed efficiente.
Nello specifico i titoli nel dataset costituiscono l'S\&P500, un indice azionario che raggruppa le 500 aziende a maggiore capitalizzazione e rilevanza economica degli Stati Uniti.
Il dataset contiene uno storico di questi titoli che va dal 2000 al 2026.
Per ogni titolo sono presenti su base giornaliera quattro valori:
\begin{enumerate}
    \item \textbf{high}: massimo prezzo raggiunto dall'asset nella giornata
    \item \textbf{low}: minimo prezzo raggiunto dall'asset nella giornata
    \item \textbf{close}: prezzo di chiusura della giornata per quell'asset
    \item \textbf{vix}: indice di paura del mercato condiviso fra tutti i titoli
\end{enumerate}

È importante specificare che il dataset è composto solo dalle aziende che fanno parte attualmente dell'S\&P500. Non compaiono quindi tutte quelle società che sono uscite dall'indice per fallimento dal 2000 a oggi.
Questo è un limite accettato in quanto l'obiettivo non è la valutazione dei rendimenti ma un confronto fra strategie che operano sullo stesso paniere.


\subsection{Feature e loro selezione}
La scelta di avere nel dataset solo questi quattro indicatori si è consolidata man mano che il modello veniva testato.
All'inizio si era partiti da un grande insieme di indicatori tecnici generati tramite la libreria di FinRL; successivamente, attraverso varie ablation, si è arrivati a questo insieme minimale.
Infatti l'ablation ha dimostrato che la decisione è dominata dalla volatilità e che la maggior parte di questa moltitudine di indicatori non aggiungeva alcuna informazione utile a stimarla con precisione più di quanto non facessero i 4 indicatori definitivi.


\subsection{Suddivisione dei dati}
La suddivisione dei dati per training, validation e test set è temporale. 
Questo per evitare che durante la simulazione di trading il modello abbia accesso a dei dati futuri rispetto al momento simulato in cui si trova.

\begin{table}[h]
\centering
\begin{tabular}{lll}
\toprule
Insieme & Periodo & Uso \\
\midrule
Training     & 2000-01-01 -- 2012-01-01 & stima dei parametri della rete \\
Validazione  & 2012-01-01 -- 2015-01-01 & selezione/controllo del modello \\
Test (out-of-sample) & 2015-01-01 -- 2026-02-01 & valutazione finale \\
\bottomrule
\end{tabular}
\caption{Suddivisione temporale del dataset.}
\label{tab:split}
\end{table}

Tutti i risultati riportati si riferiscono al periodo di test, mai
osservato durante l'addestramento. I panieri operativi per gli esperimenti sul valore della parte appresa sono composti da 10
titoli per universo. Per la parte di test sulla generalità sono stati usati panieri che si differenziano tra loro per caratteristiche e numero di asset.


% ==========================================================================
\section{Metodologia sperimentale}
\label{sec:metodi}

\subsection{Strategia appresa e formula di riferimento}
La decisione del modello è l'esposizione del capitale al mercato. 
La quota che sceglie di non tenere in liquidità viene poi ripartita equamente tra i titoli del paniere fornito in input.
L'esposizione si basa su un overlay di volatility targeting:
\[
  e_t \;=\; m_\theta \cdot \underbrace{\left(\frac{\sigma_{\text{target}}}{\sigma_{\text{mkt},t}}\right)^{p}}_{\text{volatility targeting}},
\]

dove $\sigma_{\text{mkt},t}$ è la volatilità realizzata del mercato,
$\sigma_{\text{target}}$ la volatilità obiettivo, $m_\theta$ il moltiplicatore
appreso e $p$ l'esponente. La formula fissa ($m_\theta=1$, $p$ costante) è il
riferimento non appreso rispetto al quale si misura ogni componente adattiva.

\subsection{Funzione di reward}
L'addestramento del modello tramite PPO ha l'obiettivo di massimizzare una reward che unisce rendimento e controllo del rischio. 
La reward usata \textbf{DD\_VOL\_PENALTY} è il rendimento logaritmico del portafoglio al quale vengono sottratte due penalità:
\begin{equation}
r_t = \underbrace{\ln\!\frac{V_t}{V_{t-1}}}_{\text{crescita}}
      \; - \; \underbrace{\lambda_{dd}\,\mathrm{DD}_t^{\,\alpha}}_{\text{penalità drawdown}}
      \; - \; \underbrace{\lambda_{vol}\, e_t\, \sigma_{\mathrm{mkt},t}}_{\text{penalità esposizione}\times\text{vol}} .
\end{equation}

Il primo termine è il rendimento logaritmico dello step, che viene sommato a quello precedente durante tutto l'episodio e insieme formano la crescita logaritmica totale del capitale.
Il secondo termine penalizza il drawdown dello step corrente $DD_t$ elevato a un esponente $\alpha = 1,5$. 
Il valore di $\alpha$ è stato scelto in modo da far crescere la penalità man mano che il drawdown si aggrava, dando poca urgenza alle perdite superficiali e molta a quelle profonde.
Il terzo termine penalizza l'esposizione $e_t$ in proporzione alla volatilità di mercato, rendendo esplicito il legame rischio/esposizione e facendo in modo che il modello non debba scoprire per tentativi che è da evitare una grande esposizione nei periodi di alta volatilità.
I pesi usati sono
$\lambda_{\text{dd}}=0.01$ e $\lambda_{\text{vol}}=0.002$.

\subsection{Protocollo di valutazione}
Ogni modello è valutato con un backtest statico sul periodo di test. I pesi vengono ricalcolati a ogni
passo applicando i costi di transazione. I backtest sono stati fatti su periodi che rappresentano vari regimi di mercato:

\begin{table}[h]
\centering
\begin{tabular}{lll}
\toprule
Regime & Periodo & Caratterizzazione \\
\midrule
Test completo & 2015-01-01 -- 2026-02-01 & campione out-of-sample lungo \\
2017 & 2017-01-01 -- 2017-12-31 & mercato calmo/rialzista \\
COVID & 2020-02-01 -- 2020-08-31 & crash rapido e rimbalzo \\
2022 & 2022-01-01 -- 2022-12-31 & ribasso lento e prolungato \\
2025 & 2025-02-01 -- 2026-02-01 & sell-off dei dazi (apr 2025) \\
\bottomrule
\end{tabular}
\caption{Regimi di valutazione.}
\label{tab:regimi}
\end{table}

Per ogni configurazione si addestrano più seed (42, 43, 44) e si riporta
la mediana tra i seed. Questo perché gli esperimenti preliminari hanno
mostrato che seed diversi possono dare esiti sensibilmente diversi, quindi la
robustezza al seed è parte del protocollo.

\paragraph{Confronto a parità di cash medio.}
Per verificare rigorosamente che la parte appresa del modello abbia davvero aggiunto valore è stato costruita una frontiera di riferimento. 
La frontiera è stata costruita partendo dal volatility targeting di Moreira e Muir e tracciando una curva di punti (cash medio, MDD) al variare del target di volatilità $\sigma_{target}$.
Sono state poi confrontate le performance dell'agente con la frontiera per capire se esso si è riuscito a posizionare in un'area del grafico più convenitente rispetto alla frontiera.



\subsection{Baseline di confronto}
Ogni strategia appresa è confrontata con riferimenti non appresi:
\begin{itemize}
  \item la \textbf{formula fissa} ($m_\theta=1$, $p$ costante), per isolare il
        contributo dell'apprendimento rispetto alla sola struttura;
  \item il \textbf{paniere equipesato} pienamente investito, per avere un riferimento senza gestione dell'esposizione in modo da comprendere quanto l'overlay riduce il drawdown rispetto a restare sempre esposti;
  \item la posizione \textbf{100\% cash}, come null di riferimento per il
        controllo del drawdown.
\end{itemize}

\subsection{Metriche di valutazione}
Essendo l'obiettivo la riduzione del drawdown a parità di
capacità di stare sul mercato, si adotta il seguente insieme di metriche.

\paragraph{Metriche primarie.}
\begin{itemize}
  \item \textbf{Maximum Drawdown} (MDD): massima perdita dal picco,
        $\max_t (\text{picco}_t - V_t)/\text{picco}_t$. È la metrica guida.
  \item \textbf{Calmar}: $\text{CAGR}/\text{MDD}$. Rendimento per unità di
        drawdown. Premia chi riduce il drawdown senza rinunciare al
        rendimento. Serve a capire se la difesa è ottenuta semplicemente restando in liquidità o da una strategia di esposizione intelligente.
  \item \textbf{Sharpe}: rendimento medio / volatilità (annualizzato).
        Utile in contrasto col Calmar per
        distinguere il controllo del rischio dalla sola conservatività.
  \item \textbf{Cash medio}: capitale medio tenuto in liquidità. Rende esplicito il
        costo di ogni riduzione di drawdown: un MDD più basso ottenuto
        con più cash è conservatività, non tempismo.
\end{itemize}

\paragraph{Metriche per il confronto sui panieri.}
Per valutare la selezione rispetto all'equipesato si usano inoltre:
\begin{itemize}
  \item \textbf{MDD\_gain}: riduzione di drawdown rispetto a un null a liquidità
        costante pari al cash medio del modello ($>0$ indica difesa attiva reale);
  \item \textbf{correlazione cash--drawdown}: verifica che la liquidità sia
        alzata in prossimità dei ribassi e non in modo indiscriminato;
  \item \textbf{up/down capture}: quota di rialzo catturata e di ribasso subita
        rispetto al mercato.
\end{itemize}

